% GENERATED FILE. Edit canonical lessons, then run npm run generate:books.
% canonical-source-hash: fnv1a64:1966d2e3668c74d7
% canonical-lessons: SA-C286-cauryam, SA-C286-yamikah, SA-C286-apaharati, SA-C286-saksi, SA-C286-palayate

\chapter{A Theft in the Village}
\label{ch:sa-a-theft-in-the-village}
\hlchaptermodality{286}

\begin{chapteropening}
\textbf{By the end of this chapter:} \emph{I can say what was stolen and who saw it.}

Complete the last lesson of chapter 286: I can say what was stolen and who saw it.
\end{chapteropening}

\section[\texorpdfstring{\sk{चौर्यम्} (cauryam)}{चौर्यम् (cauryam)}]{\sk{चौर्यम्} (cauryam) — a theft}

\label{lesson:SA-C286-cauryam}



\begin{quote}
Before the new one: say the Sanskrit for radio, then the Sanskrit for the internet.
\end{quote}

\subsection*{You'll want to know: \sk{चौर्यम्}}
\textbf{\sk{चौर्यम्}} — \emph{cauryam} — \textquotedblleft{}a theft\textquotedblright{}.

\begin{quote}
\sk{ह्यः} \sk{रात्रौ} \sk{ग्रामे} \sk{चौर्यं} \sk{घटितम्।}
\end{quote}

\begin{quote}
\emph{hyaḥ rātrau grāme cauryaṃ ghaṭitam.}
\end{quote}

\begin{quote}
A theft took place in the village last night.
\end{quote}

A neuter noun.

\begin{cousinweb}
\textbf{\sk{चौर्यम्}} (\emph{cauryam}) is built on \emph{caura}, \textquotedblleft{}a thief\textquotedblright{}, with the ending \emph{-ya} that makes a noun of the deed; MW gives \emph{caurya} as the act of a thief, theft (MW).
\end{cousinweb}

\begin{grammarlens}[title={What you've built}]
One more word for telling what happened.
\end{grammarlens}

\subsection*{Your turn}
\begin{itemize}
\raggedright
  \item \emph{Say it:} \emph{cauryam}
  \item \emph{Say it:} \emph{cauryam}, once more
  \item \emph{Say it:} \emph{hyaḥ rātrau grāme cauryaṃ ghaṭitam.} — A theft took place in the village last night.
  \item \emph{Recall:} say the Sanskrit for radio, then the Sanskrit for the internet, then say \emph{cauryam} again
\end{itemize}

\begin{tcolorbox}[breakable,colback=teal!4,colframe=teal!35!black,title={Before you move on}]
What does \sk{चौर्यम्} mean? (\textquotedblleft{}A theft\textquotedblright{}.) Say it once more.
\end{tcolorbox}
\section[\texorpdfstring{\sk{यामिकः} (yāmikaḥ)}{यामिकः (yāmikaḥ)}]{\sk{यामिकः} (yāmikaḥ) — a night watchman}

\label{lesson:SA-C286-yamikah}



\begin{quote}
Before the new one: say the Sanskrit for the internet, then the Sanskrit for a theft.
\end{quote}

\subsection*{You'll want to know: \sk{यामिकः}}
\textbf{\sk{यामिकः}} — \emph{yāmikaḥ} — \textquotedblleft{}a night watchman\textquotedblright{}.

\begin{quote}
\sk{ग्रामः} \sk{रात्रौ} \sk{यामिकाधीनः।}
\end{quote}

\begin{quote}
\emph{grāmaḥ rātrau yāmikādhīnaḥ.}
\end{quote}

\begin{quote}
At night the village depends on its watchman.
\end{quote}

A masculine noun.

\begin{cousinweb}
\textbf{\sk{यामिकः}} (\emph{yāmikaḥ}) is built on \emph{yāma}, \textquotedblleft{}a watch of the night\textquotedblright{} (MW: a turn of three hours), with the ending \emph{-ika}; MW gives \emph{yāmika} as a night-watcher or watchman (MW).
\end{cousinweb}

\begin{grammarlens}[title={What you've built}]
One more word for telling what happened.
\end{grammarlens}

\subsection*{Your turn}
\begin{itemize}
\raggedright
  \item \emph{Say it:} \emph{yāmikaḥ}
  \item \emph{Say it:} \emph{yāmikaḥ}, once more
  \item \emph{Say it:} \emph{grāmaḥ rātrau yāmikādhīnaḥ.} — At night the village depends on its watchman.
  \item \emph{Recall:} say the Sanskrit for the internet, then the Sanskrit for a theft, then say \emph{yāmikaḥ} again
\end{itemize}

\begin{tcolorbox}[breakable,colback=teal!4,colframe=teal!35!black,title={Before you move on}]
What does \sk{यामिकः} mean? (\textquotedblleft{}A night watchman\textquotedblright{}.) Say it once more.
\end{tcolorbox}
\section[\texorpdfstring{\sk{अपहरति} (apaharati)}{अपहरति (apaharati)}]{\sk{अपहरति} (apaharati) — makes off with; snatches away}

\label{lesson:SA-C286-apaharati}



\begin{quote}
Before the new one: say the Sanskrit for a theft, then the Sanskrit for a night watchman.
\end{quote}

\subsection*{You'll want to know: \sk{अपहरति}}
\textbf{\sk{अपहरति}} — \emph{apaharati} — \textquotedblleft{}makes off with; snatches away\textquotedblright{}.

\begin{quote}
\sk{चौरः} \sk{धनम्} \sk{अपहरति।}
\end{quote}

\begin{quote}
\emph{cauraḥ dhanam apaharati.}
\end{quote}

\begin{quote}
The thief makes off with the money.
\end{quote}

Past (he …): \sk{सः} \sk{धनम्} \sk{अपहृतवान्}, \emph{saḥ dhanam apahṛtavān}.

\begin{cousinweb}
\textbf{\sk{अपहरति}} (\emph{apaharati}) is built from the prefix \emph{apa}, \textquotedblleft{}away\textquotedblright{}, and the root √hṛ, \textquotedblleft{}to take\textquotedblright{}; MW gives \emph{apa-hṛ} as to snatch away or plunder (MW).
\end{cousinweb}

\begin{grammarlens}[title={What you've built}]
One more word for telling what happened.
\end{grammarlens}

\subsection*{Your turn}
\begin{itemize}
\raggedright
  \item \emph{Say it:} \emph{apaharati}
  \item \emph{Say it:} \emph{apaharati}, once more
  \item \emph{Say it:} \emph{cauraḥ dhanam apaharati.} — The thief makes off with the money.
  \item \emph{Recall:} say the Sanskrit for a theft, then the Sanskrit for a night watchman, then say \emph{apaharati} again
\end{itemize}

\begin{tcolorbox}[breakable,colback=teal!4,colframe=teal!35!black,title={Before you move on}]
What does \sk{अपहरति} mean? (\textquotedblleft{}Makes off with; snatches away\textquotedblright{}.) Say it once more.
\end{tcolorbox}
\section[\texorpdfstring{\sk{साक्षी} (sākṣī)}{साक्षी (sākṣī)}]{\sk{साक्षी} (sākṣī) — a witness}

\label{lesson:SA-C286-saksi}



\begin{quote}
Before the new one: say the Sanskrit for a night watchman, then the Sanskrit for makes off with.
\end{quote}

\subsection*{You'll want to know: \sk{साक्षी}}
\textbf{\sk{साक्षी}} — \emph{sākṣī} — \textquotedblleft{}a witness\textquotedblright{}.

\begin{quote}
\sk{साक्षी} \sk{चौरं} \sk{न} \sk{दृष्टवान्।}
\end{quote}

\begin{quote}
\emph{sākṣī cauraṃ na dṛṣṭavān.}
\end{quote}

\begin{quote}
The witness did not see the thief.
\end{quote}

A masculine or feminine noun, by the person meant.

\begin{cousinweb}
\textbf{\sk{साक्षी}} (\emph{sākṣī}) is the nominative of \emph{sākṣin}, which MW gives as one who sees with his own eyes, an eye-witness, and a witness in law; it sits beside MW's \emph{sākṣāt}, \textquotedblleft{}with one's own eyes\textquotedblright{}, but MW gives no derivation in the entry, so its origin is left uncertain here (MW).
\end{cousinweb}

\begin{grammarlens}[title={What you've built}]
One more word for telling what happened.
\end{grammarlens}

\subsection*{Your turn}
\begin{itemize}
\raggedright
  \item \emph{Say it:} \emph{sākṣī}
  \item \emph{Say it:} \emph{sākṣī}, once more
  \item \emph{Say it:} \emph{sākṣī cauraṃ na dṛṣṭavān.} — The witness did not see the thief.
  \item \emph{Recall:} say the Sanskrit for a night watchman, then the Sanskrit for makes off with, then say \emph{sākṣī} again
\end{itemize}

\begin{tcolorbox}[breakable,colback=teal!4,colframe=teal!35!black,title={Before you move on}]
What does \sk{साक्षी} mean? (\textquotedblleft{}A witness\textquotedblright{}.) Say it once more.
\end{tcolorbox}
\section[\texorpdfstring{\sk{पलायते} (palāyate)}{पलायते (palāyate)}]{\sk{पलायते} (palāyate) — runs away}

\label{lesson:SA-C286-palayate}



\begin{quote}
Before the new one: say the Sanskrit for makes off with, then the Sanskrit for a witness.
\end{quote}

\subsection*{You'll want to know: \sk{पलायते}}
\textbf{\sk{पलायते}} — \emph{palāyate} — \textquotedblleft{}runs away\textquotedblright{}.

\begin{quote}
\sk{चौरः} \sk{शीघ्रं} \sk{पलायते।}
\end{quote}

\begin{quote}
\emph{cauraḥ śīghraṃ palāyate.}
\end{quote}

\begin{quote}
The thief runs away fast.
\end{quote}

Past (he …): \sk{सः} \sk{शीघ्रं} \sk{पलायितवान्}, \emph{saḥ śīghraṃ palāyitavān}.

Future: \sk{पलायिष्यते}, \emph{palāyiṣyate}.

\begin{cousinweb}
\textbf{\sk{पलायते}} (\emph{palāyate}) is built on \emph{palā}, a form of \emph{parā}, \textquotedblleft{}away\textquotedblright{}, and the root √ay, \textquotedblleft{}to go\textquotedblright{}: to go away, flee (MW).
\end{cousinweb}

\begin{grammarlens}[title={What you've built}]
One more word for telling what happened.
\end{grammarlens}

\subsection*{Your turn}
\begin{itemize}
\raggedright
  \item \emph{Say it:} \emph{palāyate}
  \item \emph{Say it:} \emph{palāyate}, once more
  \item \emph{Say it:} \emph{cauraḥ śīghraṃ palāyate.} — The thief runs away fast.
  \item \emph{Recall:} say the Sanskrit for makes off with, then the Sanskrit for a witness, then say \emph{palāyate} again
\end{itemize}

\begin{tcolorbox}[breakable,colback=teal!4,colframe=teal!35!black,title={Before you move on}]
What does \sk{पलायते} mean? (\textquotedblleft{}Runs away\textquotedblright{}.) Say it once more.
\end{tcolorbox}
